\documentclass[british]{article}

\usepackage[margin=1.5in]{geometry}

\title{The dualising complex and duality theorems in complex analytic
geometry}
\author{Jean-Pierre Ramis, Gabriel Ruget}
\date{1970}

% Standard
\usepackage{amssymb,amsmath,amscd}
\usepackage{enumerate}
\usepackage{tikz-cd}
\usepackage{mathtools}
\usepackage{booktabs}


% Bibliography
\usepackage[backend=biber,maxnames=99,maxalphanames=4]{biblatex}
\addbibresource{PMIHES-38-1970-77.bib}


% Colours
\usepackage{xcolor}


% Fonts
\usepackage{mathrsfs}
\usepackage{fouriernc}


% Things for Quarto/pandoc
\usepackage{calc,array}
\usepackage{longtable}
\providecommand{\tightlist}{%
\setlength{\itemsep}{0pt}\setlength{\parskip}{0pt}}
\newcounter{none}
\usepackage{float}
\usepackage{graphicx}
\makeatletter
\newsavebox\pandoc@box
\newcommand*\pandocbounded[1]{% scales image to fit in text height/width
  \sbox\pandoc@box{#1}%
  \Gscale@div\@tempa{\textheight}{\dimexpr\ht\pandoc@box+\dp\pandoc@box\relax}%
  \Gscale@div\@tempb{\linewidth}{\wd\pandoc@box}%
  \ifdim\@tempb\p@<\@tempa\p@\let\@tempa\@tempb\fi% select the smaller of both
  \ifdim\@tempa\p@<\p@\scalebox{\@tempa}{\usebox\pandoc@box}%
  \else\usebox{\pandoc@box}%
  \fi%
}
% Set default figure placement to htbp
\def\fps@figure{htbp}
\makeatother


% Header and footer
\usepackage{fancyhdr}
\usepackage{lastpage}
\usepackage{datetime2}
\pagestyle{fancy}
\fancypagestyle{plain}{}
\fancyhf{}
\renewcommand{\headrulewidth}{0pt}
\cfoot{\small\thepage\ of \pageref*{LastPage}}
\lfoot{\footnotesize Build date: \today}


% Theorem environments
\usepackage{amsthm}
\newenvironment{itenv}[1]
  {\phantomsection\par\smallskip\noindent\textbf{#1.}\itshape}
  {\par\smallskip}
\newenvironment{rmenv}[1]
  {\phantomsection\par\smallskip\noindent\textbf{#1.}\rmfamily}
  {\par\smallskip}


% Shortcuts
\newcommand{\oldpage}[1]{\marginpar{\footnotesize$\Big\vert$ \textit{p.~#1}}}


% Hyperref
\usepackage{hyperref}
\hypersetup{colorlinks,linkcolor={purple!50!black},citecolor={purple!50!black},urlcolor={purple!80!black}}


\renewcommand{\abstractname}{Translator's note}

%% Document %%

\begin{document}

\maketitle
\thispagestyle{fancy}

\setcounter{footnote}{0}

\begin{abstract}
  This document is a translation by \href{https://thosgood.net}{Tim Hosgood} of the following:
  \begin{quote}
    J.-P.~Ramis, G.~Ruget. ``Complexe dualisant et théorèmes de dualité
    en géométrie analytique complexe''. \emph{Publications mathématiques
    de l'I.H.É.S.} (1970) pp.~77--91.
    \href{http://www.numdam.org/item?id=PMIHES_1970__38__77_0}{numdam.org/item?id=PMIHES\_1970\_\_38\_\_77\_0}
  \end{quote}
  For more information see \href{https://thosgood.net/translations}{thosgood.net/translations}.
\end{abstract}

\tableofcontents
\bigskip

%% Content %%

\providecommand{\scr}[1]{{\mathscr{#1}}}
\renewcommand{\cal}[1]{{\mathcal{#1}}}
\renewcommand{\frak}[1]{{\mathfrak{#1}}}
\renewcommand{\geq}{\geqslant}
\renewcommand{\leq}{\leqslant}

\newcommand{\sh}[1]{\scr{#1}}
\newcommand{\cat}[1]{\mathbf{#1}}

\newcommand{\OO}{\sh{O}}

\newcommand{\CC}{\mathrm{C}}
\newcommand{\HH}{\mathrm{H}}
\newcommand{\cc}{\mathrm{c}}
\newcommand{\FS}{\mathbf{FS}}
\newcommand{\DFS}{\mathbf{DFS}}
\newcommand{\QFS}{\mathbf{QFS}}
\newcommand{\QDFS}{\mathbf{QDFS}}
\newcommand{\KK}{\mathbf{K}}
\newcommand{\TT}{\mathbf{T}}

\newcommand{\Hom}{\operatorname{Hom}}
\newcommand{\shHom}{\underline{\Hom}}
\newcommand{\Ext}{\operatorname{Ext}}
\newcommand{\shExt}{\underline{\Ext}}
\newcommand{\HyperExt}{\operatorname{HyperExt}}
\renewcommand{\Im}{\operatorname{Im}}
\newcommand{\Ker}{\operatorname{Ker}}

\oldpage{77}

Let \(X\) be a \(\sigma\)-compact complex analytic manifold of
dimension~\(n\); we write \({\mathscr{O}}\) to mean the sheaf of germs
of holomorphic functions on \(X\), and \(\Omega\) to mean the sheaf of
germs of holomorphic differential forms of degree~\(n\) on \(X\). The
correspondence that, to a locally free \({\mathscr{O}}\)-module
\({\mathscr{F}}\), associates
\(\underline{\operatorname{Hom}}(X;{\mathscr{F}},\Omega)\), is
involutive. In \autocite{11}, Serre gives, for every integer~\(p\) and
every locally free \({\mathscr{O}}\)-module \({\mathscr{F}}\), a
(Fréchet quotient) topology on \(\mathrm{H}^p(X;{\mathscr{F}})\) along
with a pairing \[
  \mathrm{H}^p(X;{\mathscr{F}}) \times \mathrm{H}_\mathrm{c}^{n-p}(X;\underline{\operatorname{Hom}}(X;{\mathscr{F}},\Omega))
  \to \mathbb{C}
\] (where the subscript \(\mathrm{c}\) denotes the family of compacts of
\(X\)) which, whenever the groups \(\mathrm{H}^p(X;{\mathscr{F}})\) and
\(\mathrm{H}^{p+1}(X;{\mathscr{F}})\) are separated, makes the
(non-topologised) vector space
\(\mathrm{H}_\mathrm{c}^{n-p}(X;\underline{\operatorname{Hom}}(X;{\mathscr{F}},\Omega))\)
the topological dual of \(\mathrm{H}^p(X;{\mathscr{F}})\).

We could also mention the pairing \[
  \mathrm{H}_\mathrm{c}^p(X;{\mathscr{F}}) \times \mathrm{H}^{n-p}(X;\underline{\operatorname{Hom}}(X;{\mathscr{F}},\Omega))
  \to\mathbb{C}.
\]

Up to replacing
\(\mathrm{H}_\mathrm{c}^{n-p}(X;\underline{\operatorname{Hom}}(X;{\mathscr{F}},\Omega))\)
by \(\operatorname{Ext}_\mathrm{c}^{n-p}(X;{\mathscr{F}},\Omega)\), the
proof by Serre extends \autocite{10} to the case where \({\mathscr{F}}\)
is an arbitrary coherent \({\mathscr{O}}\)-module, but by using
difficult theorems (with the notation of Schwartz: \({\mathscr{E}}\) is
a flat \({\mathscr{O}}\)-module, and, for every \(x\in X\),
\({\mathscr{D}}'_x\) is an injective \({\mathscr{O}}_x\)-module), due to
coming from the division of distributions. We will show, following ideas
of Grothendieck that were readopted by Malgrange, that the duality
theorem for coherent sheaves on an analytic manifold is formally much
more simple than this (see also \autocite{14}), and that it can be
generalised to analytic spaces, as long as we replace \(\Omega\) by a
suitable ``dualising'' complex, and the \(\operatorname{Ext}\) by
\(\operatorname{HyperExt}\). As future work, probably very soon to
appear, we note (not to mention an article in preparation with
J.-L.~Verdier on a relative duality theorem) that Malgrange expects to
be able to use the dualising complex to give new proofs of results of
Siu \autocite{13} and Andreotti--Grauert \autocite{1}.

\section{Some lemmas on certain complexes of topological vector
spaces}\label{section-1}

We first clarify a point of vocabulary for those readers not familiar
with derived categories \autocite{7}: let \(\mathbf{A}\) be an abelian
category, and \(C(\mathbf{A})\) the category of bounded-below complexes
of objects of \(\mathbf{A}\) (with differentials of degree~\(+1\) and
morphisms of degree~\(0\)). \oldpage{78} We define \(T\) to be the
automorphism of \(C(\mathbf{A})\) that translates complexes one notch to
the left and changes the sign of all the differentials. An object of
\(C(\mathbf{A})\) is said to be \emph{acyclic} if all its cohomology
objects are zero. Let \(u\colon X^\bullet\to Y^\bullet\) be a morphism
in \(C(\mathbf{A})\); we say that it is a \emph{quasi-isomorphism} if it
induces an isomorphism of cohomology objects of \(X^\bullet\) to those
of \(Y^\bullet\); in this case, we also say that
\(\xrightarrow{u}Y^\bullet\) is a \emph{resolution} of \(X^\bullet\). We
define the \emph{cylinder} of \(u\) to be the complex
\(T(X^\bullet)\oplus Y^\bullet\) endowed with the differential
\(\left(\begin{smallmatrix}T(d_X)&0\\T(u)&d_Y\end{smallmatrix}\right)\).
Saying that \(u\) is a quasi-isomorphism is equivalent to saying that
its cylinder is zero.

We denote by the acronym \(\mathbf{FS}\) (resp. \(\mathbf{DFS}\))
Fréchet--Schwartz spaces (resp. the strong duals of Fréchet--Schwartz
spaces \autocite{6}), and by the acronym \(\mathbf{QFS}\) (resp.
\(\mathbf{QDFS}\)) the quotients of \(\mathbf{FS}\) spaces (resp. of
\(\mathbf{DFS}\) spaces) (which are, as it were, non-separated
\(\mathbf{FS}\) or \(\mathbf{DFS}\) spaces).

\protect\phantomsection\label{lemma-1}
\begin{itenv}{Lemma 1}
Let \(X^\bullet\) and \(Y^\bullet\) be complexes of Fréchet spaces (with
continuous linear differentials) and \(u\colon X^\bullet\to Y^\bullet\)
a (continuos linear) morphism. If \(u\) is a quasi-isomorphism (in the
purely algebraic sense) then it is a topological quasi-isomorphism,
i.e.~it induces isomorphisms between the cohomology spaces of
\(X^\bullet\) and of \(Y^\bullet\), where these spaces are endowed with
their natural topologies (which are not necessarily separated).

\end{itenv}

\begin{proof}
If \(H^i\) and \(K^i\) denote the cohomology spaces of \(X^\bullet\) and
\(Y^\bullet\) (respectively), then it suffices to show that \(u\) sends
the closure of \(0\) in \(H^i\) to the closure of \(0\) in \(K^i\). Let
\(M^i\) and \(N^i\) be the spaces of cycles of degree~\(i\) of
\(X^\bullet\) and \(Y^\bullet\) (respectively); these are Fréchet
spaces; thus, the map \(d_Y+u\) from \(Y^{i-1}\oplus M^i\) to \(N^i\),
which is surjective, is a homomorphism. Since the closed
\(Y^{i-1}+\overline{d_X(X^{i-1})}\) is saturated with respect to the
``projection'' \(d_Y+u\), its image is closed in \(N^i\), and is thus
\(\overline{d_Y(Y^{i-1})}\), which implies the desired property when we
pass to cohomology groups.
\end{proof}

\protect\phantomsection\label{lemma-1-bis}
\begin{itenv}{Lemma 1 bis}
\hyperref[lemma-1]{Lemma~1} remains true if we suppose that both
\(X^\bullet\) and \(Y^\bullet\) are not Fréchet complexes, but complexes
of \(\mathbf{DFS}\) spaces.

\end{itenv}

\begin{proof}
This comes from ``permanence'' properties of \(\mathbf{DFS}\) spaces,
and the closed graph theorem for these spaces \autocite{15}.
\end{proof}

\protect\phantomsection\label{lemma-2}
\begin{itenv}{Lemma 2}
Let \(E\xrightarrow{u}F\xrightarrow{v}G\) be a sequence of
\(\mathbf{FS}\) spaces and continuous linear maps, such that
\(v\circ u=0\); let \(E'\xleftarrow{u'}F'\xleftarrow{v'}G'\) be the
transpose sequence (with the duals being endowed with the strong
topology). Then the pairing between \(F\) and \(F'\) induces a pairing
between \(\operatorname{Ker}{v}/\overline{\operatorname{Im}{u}}\) and
\(\operatorname{Ker}{u'}/\overline{\operatorname{Im}{v'}}\), which makes
each of these spaces (endowed with their natural topology) the strong
dual of one another.

\end{itenv}

\begin{proof}
The natural map from
\(\operatorname{Ker}{u'}/\overline{\operatorname{Im}{v'}}\) to
\((\operatorname{Ker}{v}/\overline{\operatorname{Im}u})'\) is surjective
by Hahn--Banach, injective thanks to the semi-reflexivity of \(F\), and
continuous since \(a\) is hypocontinuous with respect to the bounded
subsets of \(\operatorname{Ker}{v}/\overline{\operatorname{Im}{u}}\) (to
see this, it suffices to know how to lift a bounded, i.e.~a relatively
compact, of \(\operatorname{Ker}{v}/\overline{\operatorname{Im}{u}}\) to
a bounded of \(\operatorname{Ker}{v}\)); it is thus bicontinuous.
\oldpage{79} Since
\(\operatorname{Ker}{v}/\overline{\operatorname{Im}{u}}\), being an
\(\mathbf{FS}\) space, is reflexive, we have thus proven the claim.
\end{proof}

\protect\phantomsection\label{lemma-3}
\begin{itenv}{Lemma 3}
Let \(X^\bullet\) and \(Y^\bullet\) be complexes of \(\mathbf{FS}\)
spaces (or even of \(\mathbf{DFS}\) spaces), and
\(u\colon X^\bullet\to Y^\bullet\) a quasi-isomorphism. Then the
transpose \(u'\) of \(u\) is a quasi-isomorphism.

\end{itenv}

\begin{proof}
The cylinder of \(u\) is acyclic; it is thus a complex whose
differentials are operators with closed image; the transpose complex,
which is the cylinder of \(u'\), also has differentials with closed
image, and, by \hyperref[lemma-2]{Lemma~2}, is acyclic.
\end{proof}

Finally, we give a lemma that allows us to cheaply ensure that the
topologies on pairs of spaces in duality coincide:

\protect\phantomsection\label{lemma-4}
\begin{itenv}{Lemma 4}
Let \(P\) and \(Q\) be \(\mathbf{QFS}\) spaces, and \(R\) and \(S\) be
\(\mathbf{QDFS}\) spaces; let \(a\) be a pairing between \(P\) and \(R\)
that puts the associated separated spaces in duality, and let \(b\) be a
pairing of the same sort between \(Q\) and \(S\); let \(u\colon P\to Q\)
and \(v\colon S\to R\) be linear maps; all such that \(u\) and \(v\) are
transpose to one another with respect to \(a\) and \(b\). Then \(u\) and
\(v\) are continuous.

\end{itenv}

\begin{proof}
Since the closure of \(\{0\}\) in \(P\) (resp. in \(S\)) is orthogonal
to \(R\) (resp. to \(Q\)), we see that
\(u\overline{\{0_P\}}\subset\overline{\{0_Q\}}\) and
\(v\overline{\{0_S\}}\subset\overline{\{0_R\}}\). The maps \(u\) and
\(v\) thus induce transpose maps \(\hat{u}\) and \(\hat{v}\) between the
separated spaces associated to \(P\), \(Q\), \(R\), and \(S\). Since
\(\hat{u}\) is a transpose, it is weakly continuous, and so its graph is
closed; it is thus strongly continuous, as is \(\hat{v}\), and the
continuity of \(u\) and \(v\) then follow.
\end{proof}

\section{Trace. Elementary duality theorem}\label{section-2}

We use the notations \(X\), \({\mathscr{O}}\), and \(\Omega\) from the
introduction. As a topological manifold, \(X\) is canonically oriented,
and Poincaré duality gives us a canonical linear form on
\(\mathrm{H}_\mathrm{c}^{2n}(X;\mathbb{C})\). This form, pre-composed
with the composition product of the \(\operatorname{Ext}\) (thinking of
the \(\operatorname{Ext}\) as sheaves of \(\mathbb{C}\)-spaces), gives a
pairing: \[
  \mathrm{H}_\mathrm{c}^n(X;\Omega) \times \operatorname{Ext}_\mathrm{c}^n(X;\Omega,\mathbb{C})
\to \mathbb{C}.
\] The complex structure of \(X\) now gives us (thinking of the
\(\operatorname{Ext}\) in terms of their Yoneda interpretation) a
remarkable \(n\)-extension of \(\Omega\) by \(\mathbb{C}\), namely: \[
  0 \to \mathbb{C} \to {\mathscr{O}}\xrightarrow{d'} \Omega^1 \xrightarrow{d'} \ldots \xrightarrow{d'} \Omega \to 0
\] where \(\Omega^i\) denotes the sheaf of holomorphic differential
forms of degree~\(i\) on \(X\). The linear form on
\(\mathrm{H}_\mathrm{c}^n(X;\Omega)\) associated to this extension is
called the \emph{trace}. The trace also arises in the following way: the
complex \[
  0 \to {{\mathscr{D}}'}^{n,0} \xrightarrow{d''} \ldots \xrightarrow{d''} {{\mathscr{D}}'}^{n,n} \to 0
\tag{1}
\] where \({{\mathscr{D}}'}^{i,j}\) denotes the sheaf of germs of
differential forms on \(X\) of type~\((i,j)\) with distribution
coefficients, is a fine resolution of \(\Omega\). The group
\(\mathrm{H}_\mathrm{c}^n(X;\Omega)\) is thus the \((n+1)\)-th
cohomology group of the complex of \(\mathbf{DFS}\) spaces \oldpage{80}
\[
  0 \to {{\mathscr{E}}'}^{n,0}(X) \xrightarrow{d''} \ldots \xrightarrow{d''} {{\mathscr{E}}'}^{n,n}(X) \to 0.
\tag{2}
\] It is thus naturally endowed with the structure of a
\(\mathbf{QDFS}\) space, and we see that integration on \(X\) of forms
of type~\((n,n)\) with compact support induces a continuous linear form
on \(\mathrm{H}_\mathrm{c}^n(X,\Omega)\), and this form is precisely
\((-1)^n\) times the trace. (Given a global section \(s\) with compact
support of \({{\mathscr{D}}'}^{n,n}\), we obtain \autocite{4} an
\(n\)-extension \(\overline{s}\), in the sense of Yoneda, of
\(\mathbb{C}\) by \(\Omega\). Pre-composed with the \(d'\)-cohomology
complex, this complex gives the extension
\((-1)^{n^2}\omega\times\overline{s}\in\operatorname{Ext}_{\mathrm{c},\mathbb{C}}^{2n}(X;\mathbb{C},\mathbb{C})\);
but we note that this is also ``equivalent'' to the \(2n\)-extension of
\(\mathbb{C}\) by \(\mathbb{C}\) associated to the de Rham complex
\[0\to\mathbb{C}\to{\mathscr{D}}'_0\to{\mathscr{D}}'_1\to\ldots\to{\mathscr{D}}'_{2n}\to0\]
and to the global section \(s\) of
\({\mathscr{D}}'_{2n}={{\mathscr{D}}'}^{n,n}\). Since Poincaré duality
is given by integration, our claim is proved.)

Now let \({\mathscr{F}}\) be an arbitrary \({\mathscr{O}}\)-module. The
composition product of the \(\operatorname{Ext}\) (of
\({\mathscr{O}}\)-modules), followed by the trace, gives us the pairings
\[
  \mathrm{H}^p(X;{\mathscr{F}}) \times \operatorname{Ext}_\mathrm{c}^{n-p}(X;{\mathscr{F}},\Omega)
  \to \mathbb{C}
\tag{3}
\] and \[
  \mathrm{H}_\mathrm{c}^p(X;{\mathscr{F}}) \times \operatorname{Ext}^{n-p}(X;{\mathscr{F}},\Omega)
  \to \mathbb{C}.
\tag{3~bis}
\]

\protect\phantomsection\label{lemma-5}
\begin{itenv}{Lemma 5}
If \(X\) is a Stein manifold, then \(\mathrm{H}_\mathrm{c}^i(X;\Omega)\)
is zero for \(i\neq n\); \(\mathrm{H}_\mathrm{c}^n(X;\Omega)\) is
separated, and the pairing~(3) puts it in duality with
\(\Gamma(X;{\mathscr{O}})\) (endowed with its usual \(\mathbf{FS}\)
structure).

\end{itenv}

\begin{proof}
Denote by \({\mathscr{E}}^{i,j}\) the sheaf of germs of differential
forms on \(X\) of type~\((i,j)\) with \(C^\infty\) coefficients. The
complex \[
  0 \to {\mathscr{E}}^{0,0} \xrightarrow{d''} \ldots \xrightarrow{d''} {\mathscr{E}}^{0,n} \to 0
\tag{1~bis}
\] is a fine resolution of \({\mathscr{O}}\); thus the groups
\(\mathrm{H}^i(X;{\mathscr{O}})\) are the cohomology groups of the
complex of \(\mathbf{FS}\) spaces \[
  0 \to {\mathscr{E}}^{0,0}(X) \xrightarrow{d''} \ldots \xrightarrow{d''} {\mathscr{E}}^{0,n}(X) \to 0
\tag{2~bis}
\] whereas the \(\mathrm{H}_\mathrm{c}^i(X;\Omega)\) are the cohomology
groups of the transpose complex~(2). The lemma then follows from
Theorem~B (which implies in particular that all the differentials of~(2)
and (2~bis) have closed image) and from \hyperref[lemma-2]{Lemma~2}.
\end{proof}

\protect\phantomsection\label{lemma-5-bis}
\begin{itenv}{Lemma 5 bis}
If \(X\) is a Stein manifold, then \(\mathrm{H}_\mathrm{c}^i(X;\Omega)\)
is zero for \(i\neq n\); the pairing~(3~bis) puts \(\Gamma(X;\Omega)\)
(endowed with its usual \(\mathbf{FS}\) structure) and
\(\mathrm{H}_\mathrm{c}^n(X;{\mathscr{O}})\) (endowed with the
\(\mathbf{DFS}\) structure obtained by considering it as a quotient of
\({{\mathscr{E}}'}^{0,n}(X)\)) in duality.

\end{itenv}

\protect\phantomsection\label{lemma-6}
\begin{itenv}{Lemma 6}
If \(X\) is a Stein manifold, and if \({\mathscr{F}}\) is an
\({\mathscr{O}}\)-module that admits a finite resolution by free
\({\mathscr{O}}\)-modules of finite type, then the groups
\(\operatorname{Ext}_\mathrm{c}^i(X;{\mathscr{F}},\Omega)\) are zero for
\(i\neq0\), \(\operatorname{Ext}_\mathrm{c}^n(X;{\mathscr{F}},\Omega)\)
is endowed with a canonical \(\mathbf{DFS}\) structure, and (3) puts it
in duality with \(\Gamma(X;{\mathscr{F}})\).

\end{itenv}

\begin{proof}
This follows from \hyperref[lemma-1-bis]{Lemma~1~bis} and
\hyperref[lemma-2]{Lemma~2}, by induction on the minimum length of the
free resolution of \({\mathscr{F}}\).
\end{proof}

\section{The dualising complex: statement of the problem. A first
duality theorem}\label{section-3}

\oldpage{81}

Given a \(\sigma\)-compact analytic space \(X\), of bounded dimension,
we want to construct a complex \(\mathbf{K}_X^\bullet\) of
\({\mathscr{O}}\)-modules that is bounded and of coherent cohomology,
along with a canonical linear form \(\mathbf{T}_X\) on
\(\mathrm{H}_\mathrm{c}^0(X;\mathbf{K}_X^\bullet)\) such that the
following is true:

\protect\phantomsection\label{theorem-1}
\begin{itenv}{Theorem }
For every coherent \({\mathscr{O}}_X\)-module \({\mathscr{F}}\) and
every integer \(p\), there exists on \(\mathrm{H}^p(X;{\mathscr{F}})\) a
unique \(\mathbf{QFS}\) structure and on
\(\operatorname{Ext}_\mathrm{c}^{-p}(X;{\mathscr{F}},\mathbf{K}_X^\bullet)\)
a unique \(\mathbf{QDFS}\) structure such that the form \(\mathbf{T}_X\)
induces a perfect pairing between the associated separated spaces.
Furthermore, the separation of \(\mathrm{H}^p(X;{\mathscr{F}})\) is
equivalent to that of
\(\operatorname{Ext}_\mathrm{c}^{n-p}(X;{\mathscr{F}},\mathbf{K}_X^\bullet)\).

\end{itenv}

If \(X\) is smooth and of dimension~\(n\), then \(\mathbf{K}_X^\bullet\)
will be a resolution of \(T^n\Omega\), and these topologies and this
pairing will be exactly those mentioned in the introduction (modulo,
obviously, the substitution of
\(\operatorname{Ext}_\mathrm{c}^{-p}(X;{\mathscr{F}},T^n\Omega)\) for
\(\operatorname{Ext}_\mathrm{c}^{n-p}(X;{\mathscr{F}},\Omega)\)).

In \hyperref[section-5]{§5} we will construct:

\begin{enumerate}
\def\labelenumi{\arabic{enumi}.}
\item
  For every manifold \(V\), a resolution \(\mathbf{K}_V^\bullet\) of
  \(T^{\dim V}\Omega_V\) such that, for every integer \(p\) and every
  \(x\in V\), \(\mathbf{K}_{V,x}^p\) is an injective
  \({\mathscr{O}}_{V,x}\)-module.
\item
  For every embedding \(f\) of one manifold \(V\) into another manifold
  \(W\), a functorial isomorphism of complexes of
  \({\mathscr{O}}_V\)-modules \[
     \overline{f}\colon \mathbf{K}_V^\bullet \xrightarrow{\approx} \underline{\operatorname{Hom}}(W;f_*{\mathscr{O}}_V,\mathbf{K}_W^\bullet)
   \] that is compatible with the traces on
  \(\mathrm{H}_\mathrm{c}^{\dim V}(V;\Omega_V)=\mathrm{H}_\mathrm{c}^0(V;\mathbf{K}_V^\bullet)\)
  and on \(\mathrm{H}_\mathrm{c}^0(W;\mathbf{K}_W^\bullet)\).
\end{enumerate}

We will use these complexes \(\mathbf{K}_V^\bullet\) to construct the
``dualising complex'' \(\mathbf{K}_X^\bullet\) of the analytic space
\(X\), as follows. Let \(U\) be an open of \(X\), realisable as an
analytic subset of a Stein manifold, and let \(\xrightarrow{\varphi}V\)
be such a realisation. The complex
\(\mathbf{K}_U^\bullet=\operatorname{Hom}(V;\varphi_*{\mathscr{O}}_U,\mathbf{K}_V^\bullet)\),
considered as an \({\mathscr{O}}_U\)-module, does not depend on the
realisation of \(U\), as we can see by covering two realisations by a
third and using the transitivity of the isomorphisms \(\overline{f}\);
it is clearly endowed with a trace
\(\mathbf{T}_U\colon\mathrm{H}_\mathrm{c}^0(U;\mathbf{K}_U^\bullet)\to\mathbb{C}\).
We obtain \(\mathbf{K}_X^\bullet\) by gluing together the
\(\mathbf{K}_U^\bullet\), and the trace \(\mathbf{T}_X\) similarly (any
two traces \(\mathbf{T}_U\) and \(\mathbf{T}_{U'}\) do indeed glue
together by Mayer--Vietoris and the vanishing of
\(\mathrm{H}_\mathrm{c}^1(U\cap U';\mathbf{K}_{U\cap U'}^\bullet)\)).

\protect\phantomsection\label{lemma-7}
\begin{itenv}{Lemma 7}
Let \(U\) be an open of \(X\), and \({\mathscr{F}}\) an
\({\mathscr{O}}_U\)-module. Suppose that \(U\) admits a realisation
\(\xrightarrow{\varphi}V\) as an analytic subset of a Stein manifold,
and that \(\varphi_*{\mathscr{F}}\) admits a finite resolution by free
\({\mathscr{O}}_V\)-modules of finite type. Then the groups
\(\operatorname{Ext}_\mathrm{c}^i(U;{\mathscr{F}},\mathbf{K}_U^\bullet)\)
are zero for \(i\neq0\), and
\(\operatorname{Ext}_\mathrm{c}^0(U;{\mathscr{F}},\mathbf{K}_U^\bullet)\)
is endowed with a canonical \(\mathbf{DFS}\) structure, and the pairing
induced by \(\mathbf{T}_U\) puts it in duality with
\(\Gamma(U;{\mathscr{F}})\).

\end{itenv}

\begin{proof}
Algebraically, the canonical morphism from
\(\operatorname{Ext}_\mathrm{c}^i(V;\varphi_*{\mathscr{F}},\mathbf{K}_V^\bullet)\)
to
\(\operatorname{Ext}_\mathrm{c}^i(U;{\mathscr{F}},\mathbf{K}_U^\bullet)\)
is an isomorphism (we can reduce to the local
\(\underline{\operatorname{Ext}}\) by a standard spectral sequence
argument; the fibres of these \(\underline{\operatorname{Ext}}\) are
equal, given the coherence of \({\mathscr{F}}\), to the
\(\operatorname{Ext}\) of the fibres; and there, the isomorphism is
evident). As for \(\Gamma(U;{\mathscr{F}})\) and
\(\Gamma(V;\varphi_*{\mathscr{F}})\), they are evidently topologically
isomorphism. \hyperref[lemma-6]{Lemma~6} then gives the topology on
\(\operatorname{Ext}_\mathrm{c}^0(U;{\mathscr{F}},\mathbf{K}_U^\bullet)\)
and proves the desired duality.
\end{proof}

\oldpage{82}

We can now prove \hyperref[theorem-1]{Theorem~1}. Let \({\mathscr{F}}\)
be a coherent \({\mathscr{O}}_X\)-module, and \(I\) a countable set of
indices. We can find a family of ``charts''
\(\varphi_i\colon\Theta_i\to\mathbb{C}^{N_i}\) and, in each
\(\mathbb{C}^{N_i}\), an open polydisc \(V_i\) such that, if we set
\(U_i=\varphi^{-1}(V_i)\), we have:

\begin{enumerate}
\def\labelenumi{\arabic{enumi}.}
\tightlist
\item
  For all \(i\), the system
  \((U_i, {\mathscr{F}}|U_i, \varphi_i|U_i,V_i)\) satisfies the
  hypotheses of \hyperref[lemma-7]{Lemma~7}.
\item
  The \(U_i\) form a locally finite (Stein) open cover of \(X\), which
  we denote by \({\mathfrak{U}}\).
\end{enumerate}

We can then apply \hyperref[lemma-7]{Lemma~7} not only to the \(U_i\),
but also to the finite intersections of the \(U_i\). Leray's~theorem,
combined with Theorem~B, tell us that the groups
\(\mathrm{H}^i(X;{\mathscr{F}})\) are the cohomology groups of the
complex of \(\mathbf{FS}\) spaces \[
  \mathrm{C}^\bullet({\mathfrak{U}},\mathrm{H}^0({\mathfrak{U}};{\mathscr{F}})).
\tag{4}
\] We have thus endowed the \(\mathrm{H}^i(X;{\mathscr{F}})\) with
\(\mathbf{QFS}\) topologies that do not depend on the cover
\({\mathfrak{U}}\), by \hyperref[lemma-1]{Lemma~1}. By
\hyperref[lemma-7]{Lemma~7}, the trace \(\mathbf{T}_X\) puts in
(topological) duality the complex of cochains in~(4) and the complex of
finite chains \[
  \mathrm{C}_\bullet^\mathrm{c}({\mathfrak{U}},\operatorname{Ext}_\mathrm{c}^0({\mathfrak{U}};{\mathscr{F}},\mathbf{K}_X^\bullet)),
\tag{5}
\] where the \(\operatorname{Ext}_\mathrm{c}\) are related to one
another by the evident arrows in the direction of increasing size
(extension by zero). To interpret the homology of~(5), we consider an
injective resolution \(L^\bullet\) of \(\mathbf{K}_X^\bullet\), and the
double complex \[
  \mathrm{C}_i^\mathrm{c}({\mathfrak{U}},\operatorname{Hom}_\mathrm{c}({\mathfrak{U}};{\mathscr{F}},L^j)).
\tag{6}
\] If we calculate the hyperhomology of~(6) by first working in the
\(j\)~direction, then we find the homology of~(5) (the spectral sequence
is degenerate by \hyperref[lemma-7]{Lemma~7}). Now we work in the
\(i\)~direction. Since \(L^j\) is injective, the sheaf
\(\underline{\operatorname{Hom}}(X;{\mathscr{F}},L^j)\) is flasque, and
thus \(\mathrm{c}\)-soft, and the second spectral sequence of~(6) is
also degenerate \autocite{3}, and the hyperhomology of~(6) reduces to
the homology of \(\operatorname{Hom}_cc(X;{\mathscr{F}},L^j)\), i.e.~to
the
\(\operatorname{Ext}_\mathrm{c}^j(X;{\mathscr{F}},\mathbf{K}_X^\bullet)\).
By \hyperref[lemma-2]{Lemma~2}, we have thus obtained \(\mathbf{QDFS}\)
topologies on the
\(\operatorname{Ext}_\mathrm{c}^p(X;{\mathscr{F}},\mathbf{K}_X^\bullet)\),
such that \(\mathbf{T}_X\) puts the associated separated spaces in
duality with the separated spaces associated to the
\(\mathrm{H}^{-p}(X;{\mathscr{F}})\) (the sign comes from the fact that
the simple complex associated to~(6) is graded by \(j-i\)).
\hyperref[lemma-4]{Lemma~4} ensures the uniqueness of these
\(\mathbf{QDFS}\) topologies.

The last claim of the theorem comes from the fact that, for the spaces
in question, a map has closed image if and only if its transpose has
closed image.

\section{A second duality theorem}\label{section-4}

\protect\phantomsection\label{theorem-2}
\begin{itenv}{Theorem 2}
Let \(X\) be a \(\sigma\)-compact analytic space of bounded dimension,
\(\mathbf{K}_X^\bullet\) its dualising complex, and \({\mathscr{F}}\) a
coherent analytic sheaf on \(X\). Then, for every integer~\(p\), the
spaces \(\mathrm{H}_\mathrm{c}^p(X;{\mathscr{F}})\) and
\(\operatorname{Ext}^{-p}(X;{\mathscr{F}},\mathbf{K}_X^\bullet)\) are
endowed with canonical topologies (\(\mathbf{QDFS}\) and
\(\mathbf{QFS}\), respectively), and the trace \(\mathbf{T}_X\) induces
a pairing between these spaces that makes the associated separated space
of one the strong dual of the associated separated space of the other.
\oldpage{83} Furthermore, the separation of
\(\mathrm{H}_\mathrm{c}^p(X;{\mathscr{F}})\) is equivalent to that of
\(\operatorname{Ext}^{1-p}(X;{\mathscr{F}},\mathbf{K}_X^\bullet)\).

\end{itenv}

Here too, by \hyperref[lemma-4]{Lemma~4}, the pair of topologies on
\(\mathrm{H}_\mathrm{c}^p(X;{\mathscr{F}})\) and
\(\operatorname{Ext}^{-p}(X;{\mathscr{F}},\mathbf{K}_X^\bullet)\) is the
only pair of \((\mathbf{QDFS},\mathbf{QFS})\) topologies such that that
\(\mathbf{T}_X\) puts the associated separated spaces in duality.

We start with a lemma analogous to \hyperref[lemma-5]{Lemma~5}.

\protect\phantomsection\label{lemma-8}
\begin{itenv}{Lemma 8}
Let

\end{itenv}


% Bibliography
\nocite{*}
\printbibliography[heading=bibintoc,title=Bibliography]

\end{document}
